\documentclass[10pt,a4paper]{amsart}
\usepackage{amsmath,amssymb,amsthm,graphicx}
\usepackage{verbatim}
\newenvironment{reference}{\comment}{\endcomment}
%\newenvironment{reference}{}{}
\newenvironment{slogan}{\comment}{\endcomment}
\newenvironment{history}{\comment}{\endcomment}

% For commutative diagrams we use Xy-pic
\usepackage[all]{xy}

% We use 2cell for 2-commutative diagrams.
\xyoption{2cell}
% Theorem environments.
%
\theoremstyle{plain}
\newtheorem{theorem}[subsection]{Theorem}
\newtheorem{proposition}[subsection]{Proposition}
\newtheorem{lemma}[subsection]{Lemma}

\theoremstyle{definition}
\newtheorem{definition}[subsection]{Definition}
\newtheorem{example}[subsection]{Example}
\newtheorem{exercise}[subsection]{Exercise}
\newtheorem{situation}[subsection]{Situation}

\theoremstyle{remark}
\newtheorem{remark}[subsection]{Remark}
\newtheorem{remarks}[subsection]{Remarks}

\numberwithin{equation}{subsection}

% Macros
%
\def\lim{\mathop{\mathrm{lim}}\nolimits}
\def\colim{\mathop{\mathrm{colim}}\nolimits}
\def\Spec{\mathop{\mathrm{Spec}}}
\def\Hom{\mathop{\mathrm{Hom}}\nolimits}
\def\Ext{\mathop{\mathrm{Ext}}\nolimits}
\def\SheafHom{\mathop{\mathcal{H}\!\mathit{om}}\nolimits}
\def\SheafExt{\mathop{\mathcal{E}\!\mathit{xt}}\nolimits}
\def\Sch{\mathit{Sch}}
\def\Mor{\mathop{\mathrm{Mor}}\nolimits}
\def\Ob{\mathop{\mathrm{Ob}}\nolimits}
\def\Sh{\mathop{\mathit{Sh}}\nolimits}
\def\NL{\mathop{N\!L}\nolimits}
\def\CH{\mathop{\mathrm{CH}}\nolimits}
\def\proetale{{pro\text{-}\acute{e}tale}}
\def\etale{{\acute{e}tale}}
\def\QCoh{\mathit{QCoh}}
\def\Ker{\mathop{\mathrm{Ker}}}
\def\Im{\mathop{\mathrm{Im}}}
\def\Coker{\mathop{\mathrm{Coker}}}
\def\Coim{\mathop{\mathrm{Coim}}}

% Boxtimes
%
\DeclareMathSymbol{\boxtimes}{\mathbin}{AMSa}{"02}

%
% Macros for moduli stacks/spaces
%
\def\QCohstack{\mathcal{QC}\!\mathit{oh}}
\def\Cohstack{\mathcal{C}\!\mathit{oh}}
\def\Spacesstack{\mathcal{S}\!\mathit{paces}}
\def\Quotfunctor{\mathrm{Quot}}
\def\Hilbfunctor{\mathrm{Hilb}}
\def\Curvesstack{\mathcal{C}\!\mathit{urves}}
\def\Polarizedstack{\mathcal{P}\!\mathit{olarized}}
\def\Complexesstack{\mathcal{C}\!\mathit{omplexes}}
% \Pic is the operator that assigns to X its picard group, usage \Pic(X)
% \Picardstack_{X/B} denotes the Picard stack of X over B
% \Picardfunctor_{X/B} denotes the Picard functor of X over B
\def\Pic{\mathop{\mathrm{Pic}}\nolimits}
\def\Picardstack{\mathcal{P}\!\mathit{ic}}
\def\Picardfunctor{\mathrm{Pic}}
\def\Deformationcategory{\mathcal{D}\!\mathit{ef}}
\usepackage{cleveref}
\crefname{originallemma}{original lemma}{original lemmas}
\Crefname{originallemma}{Original lemma}{Original lemmas}
\makeatletter
\newlabel{lemma-scheme-theoretic-image}{{082X}{source}{}{sourcepointer-1}{}}
\newlabel{lemma-scheme-theoretic-image@cref}{{[originallemma][0][]082X}{[1][0][]source}{}{}{}}
\makeatother
\title{Scheme theoretic density of an open immersion of algebraic spaces characterized by injectivity of the pushforward}
\author{The Stacks Project Authors}
\date{stacks-project:spaces-morphisms.tex:local-snapshot-20261012; GNU FDL 1.2 or later}
\begin{document}
\maketitle
\setcounter{subsection}{1}
\noindent\textit{Editorial note.} All delivered mathematical text is the author's own, copied from the pinned Stacks Project chapter file spaces-morphisms.tex; no statement or proof step was rewritten, reordered or removed.\par\medskip
\noindent\textit{Editorial note.} The proof refers to the proof of the scheme theoretic image lemma for the closure computation; that lemma is in the same chapter and the argument for the converse direction is written out in full with the affine etale neighbourhood and the element f.\par\medskip
\noindent\textit{Editorial note.} The single reference of the counted spans is recorded as a non-load-bearing navigation pointer; the argument that consumes it stays inside the counted proof span.\par\medskip
\noindent\textit{Editorial note.} Preamble: the three native excerpt spans of sources/stacks-spaces-morphisms/preamble.tex (lines 8-13, 14-18 and 154-217). The chapter file itself carries no preamble text.\par\medskip
\noindent\textit{Allowed standard prerequisites (editorial).}\begin{itemize}
\item The characterization of scheme theoretic closure by the kernel of the structure map. Lemma lemma-scheme-theoretic-image, used at proof lines 3538-3540. Turns injectivity of the restriction of O\_X -> j\_*O\_U into the closure statement.
\item Etale localization of the injectivity of O\_X -> j\_*O\_U. used at proof lines 3535-3538. Allows the statement to be tested on an arbitrary algebraic space V etale over X.
\item Affine etale neighbourhoods detect non-injectivity of a map of sheaves. used at proof lines 3543-3547. Produces the nonzero f in A vanishing on the preimage, giving the contradiction.
\end{itemize}\medskip
\ifdefined\hypertarget\hypertarget{sourcepointer-1}{}\fi
\noindent\textit{Source reference.} Original lemma 082X: focused-seeds/AGP/sources/stacks-snapshot-20261011/tags.txt, source lines 10267--10267. Navigation pointer into the Stacks Project; the standard fact it supplies is applied inside the counted proof span and the argument resting on it stays there.\par\medskip
\begin{lemma}
\label{lemma-characterize-scheme-theoretically-dense}
Let $S$ be a scheme.
Let $j : U \to X$ be an open immersion of algebraic spaces over $S$.
Then $U$ is scheme theoretically dense in $X$ if and only if
$\mathcal{O}_X \to j_*\mathcal{O}_U$ is injective.
\end{lemma}

\begin{proof}
If $\mathcal{O}_X \to j_*\mathcal{O}_U$ is injective,
then the same is true when restricted to any
algebraic space $V$ \'etale over $X$.
Hence the scheme theoretic closure of $U \times_X V$ in $V$
is equal to $V$, see proof of
Lemma \ref{lemma-scheme-theoretic-image}.
Conversely, assume the scheme theoretic
closure of $U \times_X V$ is equal to $V$ for all $V$ \'etale over $X$.
Suppose that $\mathcal{O}_X \to j_*\mathcal{O}_U$ is not injective.
Then we can find an affine, say $V = \Spec(A)$, \'etale over $X$
and a nonzero element $f \in A$ such that $f$ maps to zero in
$\Gamma(V \times_X U, \mathcal{O})$. In this case the scheme theoretic
closure of $V \times_X U$ in $V$ is clearly contained in $\Spec(A/(f))$
a contradiction.
\end{proof}

\end{document}
